\documentclass{article}
\usepackage[utf8]{inputenc}
\usepackage[dutch]{babel}
\usepackage{graphicx}
\usepackage{geometry}
\geometry{a4paper, margin=2.5cm}
\usepackage{hyperref}
\usepackage{booktabs}
\usepackage{listings}
\usepackage{xcolor}

\definecolor{codegreen}{rgb}{0,0.6,0}
\definecolor{codegray}{rgb}{0.5,0.5,0.5}
\definecolor{codepurple}{rgb}{0.58,0,0.82}
\definecolor{backcolour}{rgb}{0.95,0.95,0.92}

\lstdefinestyle{mystyle}{
    backgroundcolor=\color{backcolour},   
    commentstyle=\color{codegreen},
    keywordstyle=\color{magenta},
    numberstyle=\tiny\color{codegray},
    stringstyle=\color{codepurple},
    basicstyle=\ttfamily\footnotesize,
    breakatwhitespace=false,         
    breaklines=true,                 
    captionpos=b,                    
    keepspaces=true,                 
    numbers=left,                    
    numbersep=5pt,                  
    showspaces=false,                
    showstringspaces=false,
    showtabs=false,                  
    tabsize=2
}
\lstset{style=mystyle}

\title{Terras Sun-Seeker: Zonoverzicht voor Gentse Terrassen}
\author{Wisdom Ononiba, Yoanna Oosterlinck}
\date{Februari 2026}

\begin{document}

\maketitle

\section{Eigen API}
De API maakt het mogelijk om data-updates dynamisch en schaalbaar te beheren.\cite{web:1}\cite{web:16}
\begin{itemize}
    \item \textbf{Endpoints} voor real-time zondata per terras, met parameters zoals locatie, tijdstip en lichtintensiteit (bijv. \texttt{GET /api/sun/:lat/:lng/:time}).
    \item \textbf{Weerdata-integratie} via externe bronnen zoals Open-Meteo (geen key nodig, gratis voor non-commercial).
    \item \textbf{Event listeners} via WebSockets (Socket.io) voor automatische updates bij zonpositie- of bewolkingswijzigingen.
\end{itemize}

Setup stappen:
\begin{enumerate}
    \item \texttt{npx tsc --init; npm i express typescript @types/express @types/node ts-node nodemon mongoose socket.io @types/socket.io suncalc3}
    \item Maak \texttt{.env}: \texttt{MONGODB\_URI=mongodb://localhost:27017/terras}, \texttt{OPENMETEO\_URL=https://api.open-meteo.com}
\end{enumerate}

Voorbeeld Express router (app.ts):
\begin{lstlisting}[language=TypeScript, caption=API Router voorbeeld]
import express from 'express';
import SunCalc from 'suncalc3';
import mongoose from 'mongoose';

interface SunData { lat: number; lng: number; time: Date; intensity: number; }

const router = express.Router();

mongoose.connect(process.env.MONGODB_URI!);

router.get('/sun/:lat/:lng/:time', async (req, res) => {
  const { lat, lng, time } = req.params;
  const date = new Date(time);
  const sunPos = SunCalc.getPosition(date, +lat, +lng);
  const intensity = sunPos.altitude > 0 ? Math.sin(sunPos.altitude) * 100 : 0;
  res.json({ intensity, azimuth: sunPos.azimuth });
});

export default router;
\end{lstlisting}

\textbf{Suncalc libraries (3 alternatieven):}
\begin{itemize}
    \item Primary: suncalc3 (npm i suncalc3) -- lichtgewicht, accuraat; Pros: geen deps; Cons: basis; \url{https://www.npmjs.com/package/suncalc3}\cite{web:10}
    \item Alt1: SunCalc.js (GitHub clone) -- vanilla JS; Pros: no build; Cons: outdated; \url{https://github.com/mourner/suncalc}
    \item Alt2: NOAA Solar API -- HTTP calls; Pros: officieel; Cons: rate limits; \url{https://gml.noaa.gov/grad/solcalc/}
\end{itemize}

\textbf{Weer APIs (3 alternatieven):}
\begin{itemize}
    \item Primary: Open-Meteo (https://api.open-meteo.com/v1/forecast?latitude=51.05&longitude=3.72\&current_weather=true) -- gratis, no key; Pros: high-res; Cons: non-comm >10k/day; \url{https://open-meteo.com}\cite{web:16}
    \item Alt1: OpenWeatherMap free tier -- key nodig; Pros: veel data; Cons: 60 calls/min; \url{https://openweathermap.org/api}
    \item Alt2: Tomorrow.io free -- 5-daags forecast; Pros: hyperlocal; Cons: limits; \url{https://www.tomorrow.io/weather-api/}\cite{web:21}
\end{itemize}

\subsection{3D-modellering Data}
\subsubsection{Acculmuleren van 3D-gegevens}
Download CSV van Gent 3D blocks (1km² tiles, OBJ/GLTF): \url{https://data.stad.gent/explore/dataset/gent-in-3d/table}\cite{web:1}. Script: fetch CSV, download tiles.

Alternatieven:
\begin{itemize}
    \item Alt1: Stad.Gent API v2 -- JSON export; Pros: direct; Cons: incompleet.
    \item Alt2: OpenDataSoft -- export CSV/GeoJSON; Pros: filters; Cons: size limits; \url{https://gent.opendatasoft.com}
\end{itemize}

\subsubsection{Verwerking van 3D-gegevens}
Node script met Three.js voor shadow raycasting: \texttt{npm i three}. Cache shadows in MongoDB per uur/locatie.

\begin{lstlisting}[language=JavaScript]
const THREE = require('three');
// Load OBJ, raycast from sun position
\end{lstlisting}

\subsection{Weerdata}
\subsubsection{Accumuleren van weerdata}
Cron job (node-cron): fetch Open-Meteo elke 15min, store in MongoDB.

Alternatieven: zie API sectie.

\subsubsection{Verwerking van weerdata}
Bereken cloud factor: \texttt{cloudCover * 0.8} op intensity.

\subsection{Toeristische evenementen data}
\subsubsection{Accumuleren van toeristische evenementen data}
SPARQL/REST van \url{https://data.stad.gent/api/datasets/toeristische-evenementen-visit-gent}\cite{web:7}. Fetch top 50 events.

Alternatieven:
\begin{itemize}
    \item Alt1: VisitGent RSS; Pros: simpel; Cons: unstructured.
    \item Alt2: Eventbrite API Gent; Pros: tickets; Cons: key.
\end{itemize}

\subsubsection{Verwerking van toeristische evenementen data}
Filter events nabij terrassen, mark op map.

\subsection{Restaurant data}
\subsubsection{Accumuleren van restaurant data}
\url{https://wiki.openstreetmap.org/wiki/Overpass_API#Quick_Start_(60_seconds):_for_Developers/Programmers}

\url{https://overpass-api.de/api/interpreter}

\begin{lstlisting}[language=JavaScript]
import requests
overpass_url = "http://overpass-api.de/api/interpreter"
overpass_query = """
[out:json][timeout:25];
area[name="Gent"]->.a;
(node["amenity"="restaurant"](area.a););
out;
"""
response = requests.get(overpass_url, params={'data': overpass_query})
restaurants = response.json()
\end{lstlisting}

Alternatieven:
\begin{itemize}
    \item Alt2: Yelp API; Pros: reviews; Cons: US-focus; \url{https://yelp.com}\cite{web:23}
    \item Alt3: VisitGent terraces list; Pros: lokaal; Cons: statisch; \url{https://visit.gent.be/en/covered-outdoor-cafes-ghent}\cite{web:22}
\end{itemize}

\subsubsection{Verwerking van restaurant data}
Parse lat/lng, ratings; enrich met sun data.

\section{Client-side app die communiceert met de API: Backend-frontend communicatie}
Gebruik Express API endpoints voor data, WebSockets voor real-time updates. React frontend maakt API calls en luistert naar socket events.
\begin{itemize}
    \item API calls via Axios: \texttt{axios.get(`/api/sun/\${lat}/\${lng}\${time}}.
    \item WebSocket setup met Socket.io-client: \texttt{const socket = io(); socket.on('sunUpdate', data => ...)}.
\end{itemize}

Error handling: try-catch, retries met exponential backoff.

\section{Client-side: De React-frontend}
De React-frontend wordt gebruikt om continu data op te vragen en te visualiseren.
\begin{itemize}
    \item Interactieve kaart met Leaflet + dynamische schaduwen.
    \item \textbf{Future-Sun-Slider} component (SVG bars per uur).
    \item \textbf{Zoek- en filterfunctie} ("zonnige terrassen >80\%").
    \item Async via TanStack Query / Socket.io.
\end{itemize}

Setup: \texttt{npx create-vite@latest terras-client --template react-ts; cd terras-client; npm i leaflet react-leaflet @tanstack/react-query axios socket.io-client tailwindcss leaflet-shadow-simulator}

\textbf{Map libs (3 alt.):}
\begin{itemize}
    \item Primary: Leaflet + leaflet-shadow-simulator (npm i leaflet-shadow-simulator); Pros: licht, gratis; \url{https://github.com/ted-piotrowski/leaflet-shadow-simulator}\cite{web:25}
    \item Alt1: Mapbox GL JS -- 3D shadows; Pros: mooi; Cons: key/billing.
    \item Alt2: OpenLayers -- extensible; Pros: features; Cons: complex.
\end{itemize}

\subsection{Interactie en visualisatie}
\begin{itemize}
    \item Island 3D-moddellering - Rond Gent is het leeg
    \item Lege ruimte rond de model opvullen met legende, Future-Sun-Slider, filter-mogelijkheden, info per locatie, weerdata, zoekbalk.
    \item Zonsintensiteit via heatmap (leaflet-shadow-simulator).
    \item Schaduwrendering via leaflet-shadow-simulator.
    \item Future-Sun-Slider: SVG cirkeldiagram met kleur per intensiteit.
\end{itemize}

\subsection{Zoek- en filterfunctie}
useState voor filters, query API met params.
Zoeken enkel voor terassen, evenementen, restaurants.

Zoeken op:
\begin{itemize}
    \item Naam terras
    \item Zonscore (intensiteit)
    \item Weercondities (bv. "alleen zonnige terrassen")
    \item Evenementen nabij terras
    \item Restaurant rating
\end{itemize}

\begin{lstlisting}[language=JavaScript]
const [filter, setFilter] = useState({ sunnyOnly: false });
const { data: sunData } = useQuery(['sun', lat, lng, time, filter], fetchSun);
\end{lstlisting}

Zoeken via MongoDB queries in API, bv. \texttt{db.terrassen.find({ intensity: { \$gt: 80 } })}.
Locaties zijn server-side gefilterd op basis van intensiteit, evenementen, ratings.

\subsection{Asynchrone communicatie}
TanStack Query: \texttt{useQuery(['sun', lat, lng], fetchSun)} met polling.
WebSockets: \texttt{socket.on('sunUpdate', data => setSunData(data))} voor real-time updates.


\subsection{Shadow rendering}
Integreer ShadeMap:
\begin{lstlisting}[language=JavaScript]
import ShadeMap from 'leaflet-shadow-simulator';
const shadeMap = new ShadeMap({ date: new Date(), color: '#01112f' }).addTo(map);
\end{lstlisting}

\subsection{Future-Sun-Slider component}
SVG circle met hourly arcs, kleur per intensity.


\subsection{Weerdata-integratie}
Floating weather widget met huidige condities (UV-index, wind, temp) via Open-Meteo API.

\subsection{Kaartelementen}
Markers voor:
\begin{itemize}
    \item Terrassen: kleur op basis van zonintensiteit.
    \item Evenementen: iconen voor nabijheid.
    \item Restaurants: sterren-rating.
\end{itemize}

Popup met details bij klikken: naam, intensiteit, weerdata, evenementen.

Legende met kleurcodes, iconen, filteropties.

\subsection{Client-side optimalisaties}
\begin{itemize}
    \item Caching: TanStack Query.
    \item Refresh: polling 5min.
    \item Lazy: React.lazy voor modals.
    \item Progressieve rendering: skeletons tijdens laden.
\end{itemize}

\section{Implementatie}
\begin{itemize}
    \item \textbf{API}: Node.js + TS + Express + MongoDB.
    \item \textbf{Client}: React + TS + Leaflet.
    \item \textbf{Database}: MongoDB voor zondata, evenementen, restaurants.
    \item \textbf{Versiebeheer}: GitHub + CI/CD (GitHub Actions).
\end{itemize}

\end{document}
